% Ad Atticum 7.3 --- headnote
% ad-atticum-07-03, 9 December 50 BC, Trebula (in the villa of Pontius)

\textit{Cicero to Atticus, written from the Trebulan villa of
Pontius on the fifth day before the Ides of December 50 BC
(the manuscript dateline: \textit{Scr. in Trebulano v Id. Dec.
a. 704 (50)}). Cicero is now in the middle of the overland
journey up the Via Appia from Brundisium toward Rome, with
Philotimus having caught him at Aeculanum on the eighth and
delivered an autograph letter from Atticus that he is now
answering at length. This is the substantial letter of the
homeward cluster: twelve sections, organized as Cicero
characteristically organizes them --- the public business
first, the private last, with the seam announced
(\emph{nunc venio ad privata}).}

\textit{The public business is the constitutional crisis. The
opening sections defend his quick exit from Cilicia (he did
not stay longer because no one stays longer than the senatus
consultum allows); admit the triumph is a complication he
would happily put down if the alternative is freedom of
speech in the Senate; and confess that he cannot govern in
crisis like the figure in the sixth book of \emph{De re
publica} when the alternative is to imitate Volcatius or
Servius. The diagnosis in \S 4 is the famous one: \emph{de
sua potentia dimicant homines hoc tempore periculo civitatis}
--- ``it is for their own power that these men are fighting,
at this time, with the state itself in danger.'' He tracks
back through the failures of Pompey's consulship, his own
exile, the prorogation of Caesar's command, the law of the
ten tribunes; the consequence is that everything now rests
on one citizen, and he would rather Pompey had not built
Caesar up than now have to resist him at full strength.
\S 5 supplies the Homeric tag for what comes next: when the
moment comes to call him by name in the Senate ---
\emph{dic, M. Tulli!} --- he will, \emph{suntoma}, side with
Pompey; but he means privately to push Pompey toward concord.
The picture of Caesar's coalition (the condemned, the
disgraced, the youth, the urban rabble, the indebted, Cassius
among the tribunes) is the clearest sociological sketch
Cicero has yet given of it.}

\textit{The private business closes things out: Caelius
(whose shift Cicero regrets for Caelius's own sake), the
Lucceian property blocks, the Philotimus accounts (with the
old grievance about the residual sum in Asia), the staff
who had turned mutinous in the province and have now
straightened themselves out, Hortensius's bequests, and ---
in a connoisseur's joke --- the \emph{Piraeea} mistake.
Caught having written \emph{in Piraeea} (with both the
wrong preposition and the wrong case), he defends himself
not from Caecilius but from Terence, whose comedies were
thought elegant enough to be the work of Laelius; if the
\emph{dēmoi} count as towns then Sunium is one too. And he
ends, as always in this period, on Tiro: the hope of getting
him back on his feet rests with Curius. The letter goes off
from Pontius's villa, with the long Appian climb still
ahead.}
